% =============================================================================
% §1.7 重写草稿 v2（独立片段，未 patch 进主稿）
% 基于 section_1_7_rewrite_draft_zh.tex（v1）的诊断重写，决定如下：
%   1) 删除 v1 §1.7.1「结构保持参数化与可训练自由度」——其结构性内容复述 §1.5、
%      且违背已生效的 X5「参数化归位 §1.5」口径；其唯一属 §1.7 的内容（容量匹配/
%      REMUS 初始化）并入候选小节理由句、规格留 config 表。6 小节 → 5 小节。
%   2) 噪声预算表保留正文，但脚注瘦身、表头括注下沉；训练噪声调度去 runbook 化
%      （20/80/0.5/样本级掩码 → 记录于运行配置）。
%   3) 全节整段重写为连贯学术散文：消除「本节只…不报告…」式元叙述、
%      口径/单元/判据等内部行话、流程化罗列；三轴表的代码字段名下沉脚注。
% 不变项：所有公式与表数值、证据口径、红线声明；label 仅删 subsec:structure-preserving-param
%   （已确认无他处 \ref 引用），其余 label 沿用 v1 以保持交叉引用一致。
% 红线：三轴正交；噪声 profile 仅为 filtered-state 鲁棒性正则、非导航后验；§1.7 不
%       预锁 §1.8 证据子集；不报告性能数值/排序/「最重要结构」结论。
% =============================================================================

\section{实验设置与比较协议}
\label{sec:training-baselines-protocol}

第~\ref{sec:structured-continuous-model} 节给出的结构保持参数化，已在可训练函数类层面保证了正定逆质量、半正定耗散与斜对称零功率耦合等结构约束，第~\ref{sec:energy-power-properties} 节进一步给出了静水机械核心的功率平衡性质。这些性质由参数化构造保证、不随训练改变；而模型作为长期轨迹预测器的实际行为，则取决于数值积分、状态转换、数据分布与初值扰动的共同作用，须经统一的实验设置加以检验。本节即规定这一实验设置与比较协议：控制块训练目标（第~\ref{subsec:training-objective} 小节）确立模型在受控仿真数据上的拟合方式；长期自由递推验证协议（第~\ref{subsec:rollout-protocol} 小节）将块级拟合延伸为暴露误差累积的主要验证手段；评价指标与内部诊断量（第~\ref{subsec:metrics} 小节）界定外部任务性能与结构诊断的分工；协议轴与初值扰动设计（第~\ref{subsec:protocol-axes} 小节）组织训练与评估两阶段的鲁棒性干预；候选模型、结构消融与证据纳入规则（第~\ref{subsec:candidates-evidence} 小节）则确定结构对照的比较空间与进入定量分析的准则。性能数值、模型排序及结构作用的判定均属第~\ref{sec:results-structure-evidence} 节，由统一证据集报告；本节的任务，是使这些判定得以可复现、可比较。

\subsection{控制块训练目标}
\label{subsec:training-objective}

训练目标规定模型在受控仿真数据上的拟合方式。训练数据由控制块序列组成：在单个控制块内，控制命令与外源上下文（海流、深度）保持分段常值，模型从块初值出发积分生成预测轨迹，
\begin{equation}
  y_0=\mathcal{T}_{d\to m}(s_0),\qquad
  \hat y(t_i)=\Phi_\theta(t_i;t_0,y_0),\qquad
  \hat s_i=\mathcal{T}_{m\to d}(\hat y(t_i)),
  \label{eq:block-prediction}
\end{equation}
其中 \(\mathcal{T}_{d\to m}\) 与 \(\mathcal{T}_{m\to d}\) 为式~\eqref{eq:data-to-model-state} 与式~\eqref{eq:model-to-data-state} 给出的数据态--模型态转换，\(\Phi_\theta\) 为式~\eqref{eq:complete-augmented-vector-field} 的向量场经神经常微分方程积分得到的流映射。目标轨迹同样转换到模型态 \(y_i=\mathcal{T}_{d\to m}(s_i)\)，使有海流样本的速度监督作用于相对水速度 \(\nu_r\)（依式~\eqref{eq:velocity-contract} 的速度契约），而位置、姿态与执行器状态在两种表示中保持同一数值。据此，加权训练目标取
\begin{equation}
  \mathcal{L}
  =
  \sum_i
  \lambda_x\norm{\hat x_i-x_i}^2
  +\lambda_R d_{\SO(3)}(\hat R_i,R_i)^2
  +\lambda_\nu\norm{\hat\nu_{r,i}-\nu_{r,i}}^2
  +\lambda_a\norm{\hat u_{a,i}-u_{a,i}}^2
  +\lambda_{\SO}\mathcal{L}_{\SO(3)} ,
  \label{eq:training-loss}
\end{equation}
各权重 \(\lambda_\bullet\) 记录于运行配置。姿态误差以旋转测地距离 \(d_{\SO(3)}(\hat R,R)\) 度量，而非旋转矩阵的欧氏差，使监督直接作用于 \(\SO(3)\) 流形上的几何偏离，并对脱离旋转群的非物理偏移保持敏感；\(\mathcal{L}_{\SO(3)}\) 为旋转矩阵的正交性与行列式正则项，约束预测姿态在数值积分下保持接近旋转群。速度项以相对水速度为监督对象，与第~\ref{sec:state-current-contract} 节的速度契约一致。评估与报告时，预测轨迹经 \(\mathcal{T}_{m\to d}\) 回到数据态，终端位置、姿态与总体速度等可观测量因而仍在数据空间解释（第~\ref{subsec:metrics} 小节）。

\subsection{长期自由递推验证协议}
\label{subsec:rollout-protocol}

块级训练目标只反映局部受控初值问题的拟合质量，无法说明误差在自由递推中的累积。本文因而以长期自由递推为主要验证手段：将控制块依时间串接，首个控制块由给定初值启动，其后每个控制块以上一块的预测末状态为模型态初值，并按数据协议更新位置原点、海流速度、控制命令与深度上下文。所得轨迹不再由后续真实状态校正，单块拟合误差遂沿时间逐块传递、放大，姿态漂移、速度发散、深度越界、常微分方程求解失败与非有限预测等仅在长期显现的失败模式由此暴露。

长期递推使用受控仿真生成的 AUV 轨迹，区分无海流与有海流设置；有海流数据继续依式~\eqref{eq:velocity-contract} 维持总体速度与相对水速度的转换，使位姿几何按总体速度推进、水动力配对作用于相对水速度。评估时域取 \(10\)、\(30\) 与 \(60~\mathrm{s}\) 三层，以 \(60~\mathrm{s}\) 为本章主时域；较短时域用于刻画误差随时间的增长形态与提前失败的位置，避免单一终端时刻掩盖中途已发散的轨迹。一条轨迹若未完成目标时域，其终端误差不计入该时域的条件误差统计，而由完成率与失败原因反映其可比性。为使后续证据能够区分模型结构导致的长期发散与数值、数据侧的局部失败，失败原因分为五类：真值轨迹本身失败、模型预测发散、常微分方程求解失败、非有限预测，以及深度或速度约束违反。

\subsection{评价指标与内部诊断量}
\label{subsec:metrics}

外部任务指标在数据空间定义，刻画模型作为长期轨迹预测器的可观测性能。本章以 \(60~\mathrm{s}\) 终端位置误差的中位数与第 95 百分位（P95）为并列主指标，分别刻画典型误差与尾部风险，而不以单一标量排序——这一选择的必要性将在第~\ref{sec:results-structure-evidence} 节显现，因为部分模型在中位数与尾部上的次序并不一致。完成目标时域的轨迹另报告深度误差、旋转测地误差与总体线速度、角速度误差，并可给出对应的均值、最大值或随时域增长的误差曲线。完成率用于判断评估样本是否足够健康，仅作健全性指标而不参与精度排序：终端误差只有在完成目标时域的样本上才具可比性。

内部诊断量在模型空间定义，用于说明各结构组件是否按设计工作，但不替代数据空间的外部任务误差。相对水速度误差用于核验水动力建模变量 \(\nu_r\) 是否与第~\ref{sec:state-current-contract} 节的速度契约一致；机械能量跨度与最大能量变化用于说明机械存储在递推中是否保持良定、有限的量级；\(\SO(3)\) 行列式与正交性误差用于说明姿态矩阵在长期递推中是否保持接近旋转群。各量的聚合方式统一为：先在单个运行内对评估轨迹取统计量（如终端位置误差的中位数与 P95），再在随机种子之间对运行级统计量取均值；记号 \(N\) 表示异常排除后参与聚合的有效种子数（异常准则见第~\ref{subsec:candidates-evidence} 小节）。

\subsection{协议轴与初值扰动设计}
\label{subsec:protocol-axes}

初值扰动是本章的鲁棒性设计，分别施加于训练与评估两个阶段；二者作用于不同阶段，应作为独立干预分别处理，而非合并为单一的“协议敏感性”。为此本章沿三条正交的轴组织初值扰动，如表~\ref{tab:protocol-axes} 所示：训练协议轴刻画训练阶段是否、以及如何对块初值加噪；评估配置轴刻画评估阶段初值扰动的强度与偏置类型；评估协议轴刻画评估端生成被扰动初值的采样方式。同一种轻量初值扰动既可作为训练线，也可作为评估端的采样实现，其含义由所处的轴决定；因此不同轴上的同名扰动并不等同，须分别解释。

\begin{table}[htbp]
  \centering
  \caption{初值扰动的三条正交协议轴}
  \label{tab:protocol-axes}
  \small
  \begin{tabularx}{\textwidth}{>{\raggedright\arraybackslash}p{0.20\textwidth}>{\raggedright\arraybackslash}p{0.24\textwidth}>{\raggedright\arraybackslash}X}
    \toprule
    协议轴 & 取值 & 含义 \\
    \midrule
    训练协议轴 & 干净；逐控制块独立噪声；轻量协议变体 & 训练阶段初值扰动的实现结构：干净初值；每个控制块独立采样噪声初值；或先生成轨迹级噪声观测、再令同一轨迹各控制块共享同一次噪声实现。 \\
    评估配置轴 & 干净；名义；退化；航向偏置 & 评估阶段初值扰动的强度与偏置类型，由干净到递增强度，以及叠加固定航向偏置的有偏配置；海流设置下另可附加海流估计偏置配置。 \\
    评估协议轴 & 干净；逐样本独立；轨迹一致 & 评估端生成被扰动初值的采样方式，与训练协议轴的实现结构对应而作用于评估阶段。 \\
    \bottomrule
  \end{tabularx}
  \par\vspace{2pt}
  \footnotesize 注：三条轴在实现中分别对应配置字段 \texttt{train\_protocol}、\texttt{eval\_profile} 与 \texttt{eval\_protocol}。
\end{table}

初值扰动的噪声模型遵循“仅对初值加噪、基于 profile、在常微分方程语义上一致”的原则。噪声只作用于 \(t=0\) 的初始状态：由干净数据态转为干净模型态后，在模型态的相对水速度 \(\nu_r\)、姿态 \(R\)、执行器反馈 \(u_a\) 与海流 \(v_c\) 通道上采样零均值扰动，得到被扰动的模型态初值再启动递推。如此控制的是模型实际消费的变量误差预算，并在海流设置下依式~\eqref{eq:velocity-contract} 保持 \(R\)、\(\nu_r\) 与 \(v_c\) 的语义一致，避免数据态扰动看似不大、而模型态输入已被过度污染的失配。这一接口的定位是滤波后状态的鲁棒性正则，而非完整导航误差后验：各通道误差预算主要由 profile 常数与训练集统计量确定，目标是约束模型消费变量的误差半径，而非复现导航滤波器的联合协方差结构。

各通道的 profile 预算以 REMUS~100 的航位推算与惯性量级为参照确定，列于表~\ref{tab:noise-budget}。相对速度噪声按通道自然尺度取 \(\sigma_i=\max(\text{floor}_i,\ \alpha\,\mathrm{std}_i)\)，其中 \(\mathrm{std}_i\) 取自训练集统计、\(\alpha\) 由 profile 决定、下限 \(\text{floor}_i\) 保证噪声不至不现实地趋零；姿态误差为 yaw 主导的各向异性小角度扰动，经 \(\SO(3)\) 指数映射作用于旋转矩阵。由此，名义与退化配置的部分语义是相对训练数据分布定义的（经 \(\alpha\,\mathrm{std}_i\)），而非完全固定的绝对物理规格：同一数据集内的跨模型比较因而合理，但 profile 的实际物理强度在跨数据集时可能漂移。为此，各运行在产物中记录各通道最终落地的 \(\sigma_i\)、姿态逐轴标准差与偏置幅值，使 profile 强度可追溯，并明确其相对训练分布的定义。

\begin{table}[htbp]
  \centering
  \caption{初值扰动各通道的 profile 噪声预算（REMUS~100 航位推算/惯性量级参照）}
  \label{tab:noise-budget}
  \small
  \begin{tabularx}{\textwidth}{>{\raggedright\arraybackslash}X cccc}
    \toprule
    通道（参数） & 名义训练 & 名义评估 & 退化评估 & 航向偏置评估 \\
    \midrule
    相对线速度 \(\nu_r\)（\(\alpha\)）& 0.03 & 0.05 & 0.10 & 同名义 \\
    相对角速度 \(\nu_r\)（\(\alpha\)）& 0.03 & 0.05 & 0.10 & 同名义 \\
    姿态（yaw 主导小角 / rad）& 0.0035 & 0.0050 & 0.0120 & 名义 \(+\) 固定 0.015 偏置 \\
    海流 \(v_{cx},v_{cy}\) / \(v_{cz}\)（m/s）& 0.008 / 0.004 & 0.012 / 0.006 & 0.030 / 0.015 & 同名义 \\
    执行器 \(\delta_r,\delta_s\)（rad）/ 转速（rpm）& 0.002 / 3 & 0.003 / 5 & 0.008 / 15 & 同名义 \\
    \bottomrule
  \end{tabularx}
  \par\vspace{2pt}
  \footnotesize 注：相对速度以 \(\sigma_i=\max(\text{floor}_i,\alpha\,\mathrm{std}_i)\) 取值，表列 \(\alpha\)；线速度与角速度下限分别为 \(0.005~\mathrm{m/s}\) 与 \(0.0015~\mathrm{rad/s}\)。航向偏置评估在名义评估的随机扰动上叠加固定 \(0.015~\mathrm{rad}\) 的 yaw 偏置，符号按样本固定、一条递推内不变。
\end{table}

训练端对噪声强度采用预热--爬坡--混合的渐进调度，以避免噪声初值在优化初期破坏训练稳定性：初期以干净初值优化，随后将噪声强度线性提升至目标值，并在稳态阶段令部分训练样本使用噪声初值；预热轮数、爬坡区间与混合比例记录于运行配置。评估端则对所有模型固定噪声随机种子，使不同模型之间的被扰动初值可直接比较。

\subsection{候选模型、结构消融与证据纳入规则}
\label{subsec:candidates-evidence}

\paragraph{候选比较空间。}候选比较空间由本文的可训练模型族确定，用于在同一状态契约、训练损失与长期递推评估协议下形成结构对照，而非预设性能排序。比较链条覆盖完整结构化模型、端口哈密顿风格基线、\(\SE(3)\) 加速度与动量黑箱、全状态黑箱，以及完整模型内部的结构消融，对应显式位姿几何、动量与质量映射、机械存储与广义力组织、耗散与零功率耦合四类检验维度。为使结构对照不被网络容量差异混淆，各候选在容量匹配的网络宽度下比较，海流任务统一以 REMUS~100 的质量矩阵与执行器先验初始化（见表~\ref{tab:experimental-config}），从而使模型族内的差异可在容量可比的前提下归因于结构配置而非参数规模。其中 \(\SE(3)\) 几何与动量基线、端口哈密顿风格基线参照 \(\SE(3)\) 哈密顿及李群端口哈密顿神经常微分方程的建模思路~\citep{duong2021se3hamnode,duong2024liegroupphnode}，但均在统一的 AUV 状态契约、数据协议与评估口径下重新实现为受控候选，而非直接复刻文献中的原始模型与训练设置。表~\ref{tab:model-comparison-roles} 汇总各模型类别在比较空间中的角色。

\begin{table}[htbp]
  \centering
  \caption{候选模型比较与结构消融的检验目标}
  \label{tab:model-comparison-roles}
  \begin{tabularx}{\textwidth}{>{\raggedright\arraybackslash}p{0.26\textwidth}>{\raggedright\arraybackslash}p{0.31\textwidth}>{\raggedright\arraybackslash}X}
    \toprule
    模型类别 & 代表模型 & 检验目标 \\
    \midrule
    完整结构化模型 & AUVHamNODE & 同时包含速度契约、\(\SE(3)\) 运动学、机械存储、\(D_\theta\)/\(J_\theta\)/\(\tau_\theta\) 分解和执行器滞后，作为本文结构化假设空间的完整候选项。 \\
    端口哈密顿风格基线 & 合并非保守广义力基线、位形广义力基线 & 比较机械存储核心保留时，不同非保守广义力组织方式对长期递推行为的影响。 \\
    \(\SE(3)\) 加速度黑箱 & 显式 \(\SE(3)\) 运动学与黑箱加速度动力学 & 检验保留位姿几何但不保留完整机械存储分解时的候选对照。 \\
    \(\SE(3)\) 动量黑箱 & 显式 \(\SE(3)\) 运动学、常质量映射与黑箱动量动力学 & 检验动量坐标和质量映射保留时，黑箱动量动力学的候选对照。 \\
    全状态黑箱 & 非结构化状态导数模型 & 提供缺少显式位姿几何和机械能量结构时的长期递推风险对照。 \\
    结构消融 & 无质量先验、对角耗散、无零功率升力耦合、广义力仅执行器条件化 & 提供质量先验、耦合耗散、斜对称零功率耦合和开放广义力条件化范围的逐项结构对照。 \\
    \bottomrule
  \end{tabularx}
\end{table}

\paragraph{结构消融。}结构消融在完整模型上逐项移除或弱化某一结构组件，用于评估第~\ref{sec:energy-power-properties} 节中各结构配置对长期递推行为的影响。每个消融对象对应一个明确的候选结构假设：

\begin{itemize}
  \item 质量先验消融移除或弱化基于物理质量矩阵的先验初始化，同时保留正定质量参数化，用于评估质量先验对动量--速度配对学习和长期递推行为的影响。
  \item 对角阻尼消融将半正定耦合耗散矩阵 \(D_\theta(\xi)\) 限制为非负对角形式，用于评估跨通道耗散耦合相对独立通道耗散的作用。
  \item 零功率耦合消融移除斜对称分支 \(J_\theta(\xi)\nu_r\)，用于评估升力或横向水动力类零功率耦合在候选模型族内的作用；该消融不等同于证明真实水动力升力项被唯一辨识。
  \item 广义力条件化消融把广义力分支 \(\tau_\theta\) 收缩为主要由执行器状态或命令驱动，弱化相对速度和海流上下文条件化，用于评估开放广义力通道的条件化范围对长期递推的影响。
\end{itemize}

\paragraph{统一实验配置。}为保证长期递推结果可复现，所有候选模型在统一的数据、训练与数值实现配置下比较\footnote{连续时间积分基于 torchdiffeq~\citep{chen2018neuralode} 的固定步长四阶 Runge--Kutta 实现。}，命令级参数、噪声预算与运行元数据完整保留在各运行的配置、评估摘要与溯源记录中。表~\ref{tab:experimental-config} 汇总与论证相关的主要设置。

\begin{table}[htbp]
  \centering
  \caption{长期递推实验的主要配置}
  \label{tab:experimental-config}
  \begin{tabularx}{\textwidth}{>{\raggedright\arraybackslash}p{0.18\textwidth}>{\raggedright\arraybackslash}X}
    \toprule
    配置维度 & 设置 \\
    \midrule
    数据规模 & 海流数据集含一千条轨迹、每条一百五十个控制块；控制块时长 0.2~s，状态采样间隔 0.05~s，训练与验证按 \(0.8\) 划分。 \\
    速度契约 & 数据态保存总体速度，训练前按式~\eqref{eq:velocity-contract} 转换为模型态相对水速度。 \\
    数值积分 & 固定步长四阶 Runge--Kutta，单精度浮点状态；训练、块级评估与长期递推共用同一求解器与损失口径。 \\
    网络容量 & 完整结构模型、结构消融与黑箱基线采用容量匹配的网络宽度；海流任务以 REMUS 质量矩阵与执行器先验初始化。 \\
    优化 & AdamW，学习率经线性预热后余弦衰减，至多三百个 epoch；批量、学习率端点与正则权重记录于运行配置。 \\
    递推评估 & 重采样初值，覆盖 PRBS、CHIRP 与 OU 三类控制场景，每场景 30 条评估轨迹；评估时域取 10、30 与 60~s，主时域 60~s。 \\
    \bottomrule
  \end{tabularx}
\end{table}

\paragraph{证据纳入与异常准则。}进入定量比较的运行须满足统一的数据协议、状态转换与评估协议，并完整记录随机种子、程序版本与运行环境。本章对运行级异常区分两类相互独立、触发后均移出定量聚合并单独披露的情形。其一为训练异常：训练阶段出现灾难性梯度失败（如训练日志出现“无成功训练批次”特征且计数大于零）。其二为递推发散：训练收敛，但长期递推的关键终端指标缺失、非有限或超过预设发散阈值。两类异常以训练异常优先计入，同一运行不重复计数；移出定量聚合后按类别与随机种子单独披露，以免选择性举证。若某模型在同一训练--评估协议组合下全部种子均递推发散，则不报告有限中位数，而记为长期稳定性失败。当不同代码环境或数据镜像的证据相互冲突时，以数据协议一致性、记录完整性与失败可复现性判定统一证据状态，而非按来源一概取舍，其完整判定规则见溯源与数据目录文档。至于实际进入定量比较的证据子集——即哪些 \((\text{模型}\times\text{协议}\times\text{种子})\) 组合在统一代码环境与数据镜像下被重建并纳入——则由第~\ref{sec:results-structure-evidence} 节在当前证据集口径下报告，本节不预先锁定。
